\documentclass[spanish,11pt,a4paper]{article}
\usepackage[margin=1in]{geometry}
\usepackage{fontspec}
\setmainfont{lmroman10-regular.otf}[BoldFont=lmroman10-bold.otf,ItalicFont=lmroman10-italic.otf,BoldItalicFont=lmroman10-bolditalic.otf]
\setsansfont{lmsans10-regular.otf}[BoldFont=lmsans10-bold.otf]
\setmonofont{lmmono10-regular.otf}[Scale=MatchLowercase]
\newfontfamily\LeanSymbols{seguisym.ttf}[Path=typeset_runtime/]
\newlength{\LeanCharWidth}
\usepackage{amsmath,amssymb}
\usepackage[spanish,es-noshorthands,es-tabla]{babel}
\usepackage{microtype,parskip}
\usepackage{graphicx,float,caption}
\usepackage{longtable,booktabs,array,calc}
\usepackage{xcolor}
\usepackage{xurl}
\usepackage{hyperref,bookmark}
\usepackage{etoolbox,needspace}
\setcounter{secnumdepth}{-\maxdimen}
\setlength{\emergencystretch}{3em}
\providecommand{\tightlist}{\setlength{\itemsep}{0pt}\setlength{\parskip}{0pt}}
\newcounter{none}
\AtBeginEnvironment{longtable}{\Needspace{8\baselineskip}\small}
\pretocmd{\section}{\Needspace{5\baselineskip}}{}{}
\pretocmd{\subsection}{\Needspace{4\baselineskip}}{}{}
\pretocmd{\subsubsection}{\Needspace{4\baselineskip}}{}{}
\captionsetup{font=small,labelfont=bf}
\urlstyle{same}
\providecommand{\pandocbounded}[1]{#1}
\hypersetup{pdftitle={Paper II: formalización y reproducibilidad},pdfauthor={Juan Pablo Traverso Gianini},pdflang={es-ES},colorlinks=true,linkcolor=blue,urlcolor=blue,citecolor=blue}
\title{Paper II: formalización y reproducibilidad}
\author{Juan Pablo Traverso Gianini\\\small Investigador independiente, Santiago, Chile\\\small \href{mailto:jtraverso@gmail.com}{jtraverso@gmail.com}\\\small ORCID: 0009-0003-6068-4096}
\date{}
\begin{document}
\raggedbottom
\maketitle
\textbf{Paper II · Formal companion v1.3}\\
19 de octubre de 2026

Este documento acompaña al artículo matemático. Conserva las
declaraciones formales, las precisiones de alcance y la evidencia
computacional que antes figuraban en el artículo. El congelado de
fuentes y su verificación registrada no cambian. No se ejecutó Lean/Lake
para esta edición. Los números históricos de secciones y tablas que
siguen remiten al artículo anterior; se conservan como localizadores y
no se renumeran como resultados nuevos de verificación. Las referencias
originales {[}1{]}, {[}2{]}, etc. figuran al final.

\href{https://github.com/jtraverso/erdos-81-chordal-clique-partitions/tree/eb335845bfa23375d5040ea31fc4f4b366346287/preprints/PAPER_II/05_formalization}{Frozen
sources and historical evidence / Fuentes congeladas y evidencias
históricas} ·
\href{https://github.com/jtraverso/erdos-81-chordal-clique-partitions/tree/eb335845bfa23375d5040ea31fc4f4b366346287/preprints/PAPER_II/01_manuscript}{Previous
article / Artículo anterior}

\subsubsection{\texorpdfstring{Corolario 1.2\(^{\prime}\) (asintótica
afinada)}{Corolario 1.2\^{}\{\textbackslash prime\} (asintótica afinada)}}\label{corolario-1.2prime-asintuxf3tica-afinada}

\textbf{Estado de la formalización.} El freeze Paper II Lean v1.2
entregado usa Lean 4 / Mathlib v4.28.0 {[}5,6{]}. Su build principal
registrado completó 8.061 jobs, y el suplemento explícito de
extremizadores y defectos de copia completó 8.032 jobs. La declaración
principal se llama \nolinkurl{PaperII.theorem_1_2}; corresponde al
Teorema 1.1 de este manuscrito e incluye el enunciado de existencia del
maximizante que lo acompaña. El mismo freeze incluye el argmax
complete-split, los conjuntos de nivel, los defectos de copia y
\nolinkurl{PaperII.AsymptoticCorollaries}. Su gate de axiomas a nivel de
teorema no registra axiomas matemáticos propios del proyecto. El
desarrollo formal deriva la estructura cordal usada en la demostración a
partir del predicado estándar \nolinkurl{IsChordal}; no se añade ningún
axioma sobre árboles de cliques ni sobre estructura cordal. La Sección 8
registra el perímetro exacto del freeze.

\subsection{8. Reproducibilidad, verificación formal y
alcance}\label{reproducibilidad-verificaciuxf3n-formal-y-alcance}

La demostración del artículo es analítica. Durante el desarrollo se
usaron experimentos computacionales para buscar contraejemplos y
detectar errores de transcripción, pero ningún cálculo se usa como
premisa de la demostración.

El ledger de cálculos acompañante registra la demostración con un nivel
de granularidad menor. Enumera las definiciones, la construcción de
cobertura copiada de Lema 3.1, las dos identidades de cancelación, la
elevación por convexidad discreta a clases de clones, el argumento de
progreso de clones, el requisito de simplicialidad en la reducción
cordal, la caracterización terminal complete-split, el programa de dos
variables sobre \(S_{p,q}\), el caso degenerado \(S_{2,0}\), las
fórmulas para \(\Phi_\tau(S_{p,q})\) y la maximización entera final. El
teorema, la cota exacta y el ensamblaje han recibido una auditoría
adversarial interna. Esto registra el proceso de validación del autor y
no es una afirmación de revisión externa por pares.

\subsubsection{Verificación formal en
Lean}\label{verificaciuxf3n-formal-en-lean}

El snapshot consolidado en
\nolinkurl{PAPER_II/05_formalization/lean_v1.2_freeze} usa Lean 4 /
Mathlib v4.28.0 {[}5,6{]}. El build principal registrado completó 8.061
jobs sin errores; un suplemento explícito que cubre
\nolinkurl{PaperII.Extremizer} y \nolinkurl{PaperII.CopyDefect} completó 8.032
jobs sin errores. La declaración principal correspondiente es

\[
\texttt{PaperII.theorem\_1\_2}
\]

en el repositorio {[}7{]}. El teorema formal incluye ambos lados del
resultado: la cota superior sobre todos los grafos cordales de \(n\)
vértices y la existencia de un grafo complete-split que alcanza el
valor.

La formalización usa el predicado cordal estándar \nolinkurl{IsChordal}. En
el desarrollo, este predicado dice que todo ciclo de longitud al menos
cuatro tiene una cuerda. El funcional de déficit se formaliza como

\[
\text{\LeanSymbols phiTau G = (|E(G)| : ℝ) - 2 · tau3star G},
\]

donde \nolinkurl{tau3star} es el número fraccional de cobertura de
triángulos del PL finito de coberturas. La familia complete-split se
formaliza como \nolinkurl{completeSplit\ p\ q}, y el valor extremo es el
entero \(\lfloor(2n+1)^2/24\rfloor\).

No se añade ningún axioma estructural cordal. La formalización deriva de
\nolinkurl{IsChordal} la propiedad hereditaria para subgrafos inducidos, la
preservación al añadir un vértice con vecindario clique y la existencia
de vértices simpliciales. La demostración escrita de Lema 5.1 usa
árboles de cliques como infraestructura estándar de libro de texto; la
demostración Lean usa un argumento de caracterización terminal sin
árboles de cliques. Por tanto, la demostración formal no supone un
axioma de árbol de cliques ni de intersección consecutiva.

El desarrollo comprobado cubre la desigualdad de cobertura copiada, la
convexidad discreta de clases de clones, el progreso de las clases de
clones, la copia que preserva la cordalidad, la caracterización terminal
complete-split, el promedio por órbitas sobre \(S_{p,q}\), el programa
lineal terminal y la maximización entera exacta. El log del freeze
registra 49 advertencias del linter, principalmente variables de sección
y argumentos \nolinkurl{simp} no utilizados; ninguna afecta la elaboración
por el kernel. El informe de axiomas registrado es

\begin{quote}
{\LeanSymbols\fontsize{8}{10}\selectfont
\noindent \nolinkurl{#print}\hspace{0.45em}\nolinkurl{axioms}\hspace{0.45em}\nolinkurl{PaperII.theorem_1_2}\par
\noindent \nolinkurl{[propext,}\hspace{0.45em}\nolinkurl{Classical.choice,}\hspace{0.45em}\nolinkurl{Quot.sound]}\par
}
\end{quote}

Estos son los axiomas clásicos estándar de Lean. El informe no contiene
\nolinkurl{sorryAx} ni axiomas propios del proyecto. El archivo de fuentes
congelado es \nolinkurl{PAPER_II_lean_v1.2_freeze.zip}, con SHA-256
\nolinkurl{ee2d05cc40d943ca92f8f7bf3e5dd83c2692518ddea5e2ca4f7686ccb1ac3895}.
Los logs principal y suplementario se almacenan como
\nolinkurl{gate_logs/BUILD_LOG.txt} y
\nolinkurl{gate_logs/BUILD_SUPPLEMENT_EXTREMIZER_COPYDEFECT.txt} dentro
de \nolinkurl{lean_v1.2_freeze}. Los módulos de contribución sobre grafos
cordales entregados en el mismo archivo son candidatos para su
incorporación upstream a Mathlib y no constituyen supuestos adicionales
de la demostración de este artículo. El archivo congelado superó una
reproducción independiente en clean room; la auditoría externa registró
un build exitoso de 8.063 jobs y las 16 superficies de axiomas
esperadas.

\textbf{Tabla 4. Perímetro de la formalización congelada (Lean v1.2).}
El freeze entregado contiene el teorema principal y las declaraciones
del manuscrito que se enumeran a continuación; la reproducción
independiente en clean room obtuvo PASS.

{\def\LTcaptype{none} % do not increment counter
\begin{longtable}[]{@{}
  >{\raggedright\arraybackslash}p{(\linewidth - 4\tabcolsep) * \real{0.3333}}
  >{\raggedright\arraybackslash}p{(\linewidth - 4\tabcolsep) * \real{0.3333}}
  >{\raggedright\arraybackslash}p{(\linewidth - 4\tabcolsep) * \real{0.3333}}@{}}
\toprule\noalign{}
\begin{minipage}[b]{\linewidth}\raggedright
Nodo
\end{minipage} & \begin{minipage}[b]{\linewidth}\raggedright
Enunciado
\end{minipage} & \begin{minipage}[b]{\linewidth}\raggedright
Huella axiomática
\end{minipage} \\
\midrule\noalign{}
\endhead
\bottomrule\noalign{}
\endlastfoot
\nolinkurl{PaperII.theorem_1_2} &
\(\max_{\text{cordal},\,|V|=n}\Phi_\tau=\lfloor(2n+1)^2/24\rfloor\),
alcanzado por un grafo complete-split & \nolinkurl{propext},
\nolinkurl{Classical.choice}, \nolinkurl{Quot.sound} \\
\nolinkurl{Fsat_argmax_unique}, \nolinkurl{Fsat_argmax_tie} & unicidad del
entero maximizante \(p\) dentro de la familia complete-split y empate
entre dos valores cuando \(n\equiv1\pmod3\) & \nolinkurl{propext},
\nolinkurl{Classical.choice}, \nolinkurl{Quot.sound} \\
\nolinkurl{level_set_iff} (Cor. 7.2) & equivalencia exacta para los
conjuntos de nivel desde el máximo continuo & \nolinkurl{propext},
\nolinkurl{Classical.choice}, \nolinkurl{Quot.sound} \\
\nolinkurl{copyDefect_nonneg}, \nolinkurl{copyGamma_ge_half_copyDefect} &
defecto de copia \(\Delta_{\Phi_\tau}\ge0\) y
\(\Gamma_{\Phi_\tau}\ge\Delta_{\Phi_\tau}/2\) & \nolinkurl{propext},
\nolinkurl{Classical.choice}, \nolinkurl{Quot.sound} \\
\end{longtable}
}

\textbf{Tabla 5. Superficie aritmética incluida en el freeze Lean v1.2.}
Estas declaraciones son consecuencias del valor extremal congelado y
están incluidas en el archivo entregado.

{\def\LTcaptype{none} % do not increment counter
\begin{longtable}[]{@{}
  >{\raggedright\arraybackslash}p{(\linewidth - 4\tabcolsep) * \real{0.3333}}
  >{\raggedright\arraybackslash}p{(\linewidth - 4\tabcolsep) * \real{0.3333}}
  >{\raggedright\arraybackslash}p{(\linewidth - 4\tabcolsep) * \real{0.3333}}@{}}
\toprule\noalign{}
\begin{minipage}[b]{\linewidth}\raggedright
Nodo
\end{minipage} & \begin{minipage}[b]{\linewidth}\raggedright
Enunciado
\end{minipage} & \begin{minipage}[b]{\linewidth}\raggedright
Huella axiomática
\end{minipage} \\
\midrule\noalign{}
\endhead
\bottomrule\noalign{}
\endlastfoot
\nolinkurl{phiTau_max_sandwich} (Cor. 1.2\(^{\prime}\)) & para todo
\(n\in\mathbb Z\),
\(4n^2+4n-23<24\lfloor(2n+1)^2/24\rfloor\le4n^2+4n+1\) &
\nolinkurl{propext}, \nolinkurl{Classical.choice}, \nolinkurl{Quot.sound} \\
\nolinkurl{odd_sq_emod_24}, \nolinkurl{phiTau_max_closed} (Cor. 7.3) &
\((2n+1)^2\equiv1\) o \(9\pmod{24}\), y la fórmula resultante por clases
de residuos para el extremo & \nolinkurl{propext},
\nolinkurl{Classical.choice}, \nolinkurl{Quot.sound} \\
\nolinkurl{phiTau_max_le_paperI_bound} &
\(\lfloor(2n+1)^2/24\rfloor\le n^2/6+n/2\) para \(n\ge1\) (comparación
con Paper I) & \nolinkurl{propext}, \nolinkurl{Quot.sound} \\
\end{longtable}
}

Las cuatro declaraciones de la Tabla 5 se reúnen en
\nolinkurl{PaperII/AsymptoticCorollaries.lean}, que compila con el mismo
toolchain (Lean 4 / Mathlib v4.28.0) y usa únicamente los axiomas
clásicos estándar de Lean, sin \nolinkurl{sorryAx} ni axiomas propios del
proyecto. Los informes \nolinkurl{#print\ axioms} de cada declaración
están archivados en \nolinkurl{PaperII/AxiomCheckCorollaries.lean}. El
módulo no añade hipótesis ni modifica el desarrollo principal congelado;
cada declaración se deriva del valor cerrado mediante \nolinkurl{ring}, una
partición finita por residuos y \nolinkurl{omega}.

\subsubsection{Verificación computacional
independiente}\label{verificaciuxf3n-computacional-independiente}

Además de la demostración analítica y del desarrollo Lean, scripts
suplementarios recalcularon las afirmaciones principales desde primeros
principios. Estos scripts evaluaron \(\tau_3^*\) mediante programación
lineal exacta, comprobaron por separado el PL dual de empaquetamiento
fraccional como prueba de consistencia, realizaron la optimización
entera con aritmética racional exacta y usaron enumeración exhaustiva
cuando fue posible. El informe de auditoría suministrado no registra
discrepancias.

Las comprobaciones reportadas incluyen:

\begin{itemize}
\tightlist
\item
  la desigualdad de copia de vértices y sus dos identidades de
  cancelación sobre todos los grafos no isomorfos de hasta \(7\)
  vértices y todos los pares no adyacentes;
\item
  convexidad discreta, progreso de clases de clones sin divisiones y
  preservación de la cordalidad, incluida una prueba sobre \(C_4\) que
  muestra por qué es necesaria la hipótesis de simplicialidad;
\item
  la caracterización terminal sobre todos los grafos cordales de hasta
  \(7\) vértices;
\item
  las fórmulas complete-split de Proposición 6.2 y Corolario 6.3
  comparadas con el PL completo en el rango ensayado;
\item
  la identidad
  \(\max_{p+q=n}\Phi_\tau(S_{p,q})=\lfloor(2n+1)^2/24\rfloor\) para
  \(1\le n\le2000\);
\item
  el máximo global sobre grafos cordales, exhaustivamente para
  \(n\le7\), con pruebas aleatorias adicionales para \(8\le n\le12\).
\end{itemize}

Los scripts de auditoría, los informes y los manifests SHA-256 acompañan
al preprint. Corroboran la demostración y ayudan a hacer reproducible el
release; no sustituyen la demostración.

\subsubsection{Alcance}\label{alcance}

El resultado probado aquí es el extremo finito exacto del funcional
fraccional \(\Phi_\tau\). No es un teorema de partición en cliques y no
usa empaquetamiento integral de triángulos. Dentro del programa Erdős
Problem \#81, el papel de este artículo es identificar la familia
terminal complete-split y el valor fraccional exacto.

\begin{center}\rule{0.5\linewidth}{0.5pt}\end{center}

\subsubsection{Corolario 7.3 (fórmula por clases de
residuos)}\label{corolario-7.3-fuxf3rmula-por-clases-de-residuos}

Así, el valor extremal es uno de estos dos enteros según la clase de
\(n\) módulo \(3\). Esto hace explícito el hecho sobre cuadrados impares
módulo \(24\) usado en el análisis por casos de la Proposición 7.1. La
misma clase residual \(n\equiv1\pmod3\) es precisamente el caso de
empate entre dos valores registrado por \nolinkurl{Fsat_argmax_tie}. El
hecho residual está formalizado como \nolinkurl{PaperII.odd_sq_emod_24}
y la forma cerrada como \nolinkurl{PaperII.phiTau_max_closed} (huella
axiomática \nolinkurl{propext}, \nolinkurl{Classical.choice},
\nolinkurl{Quot.sound}).

\subsubsection{Lema 5.1}\label{lema-5.1}

\emph{La formalización Lean acompañante prueba este lema mediante un
argumento equivalente que no usa árboles de cliques --- por eliminación
simplicial y el hecho de que todo separador minimal de vértices de un
grafo cordal es una clique --- en lugar de la demostración mediante
árboles de cliques dada abajo; véase §8. El enunciado probado es
idéntico.}

\subsubsection{Teorema 1.1}\label{teorema-1.1}

donde \(p\) es un entero más próximo a \((2n+1)/6\). Entre los grafos
complete-split \(S_{p,n-p}\), el entero maximizante \(p\) es único salvo
cuando \(n\equiv1\pmod3\); en ese caso, exactamente dos valores
consecutivos de \(p\) alcanzan el máximo. Véanse el análisis por casos
en la demostración de la Proposición 7.1 y las declaraciones
formalizadas \nolinkurl{Fsat_argmax_unique} y \nolinkurl{Fsat_argmax_tie}.
No se afirma unicidad alguna para extremizadores cordales arbitrarios.

\subsubsection{\texorpdfstring{Corolario 1.2\(^{\prime}\) (asintótica
afinada)}{Corolario 1.2\^{}\{\textbackslash prime\} (asintótica afinada)}}\label{corolario-1.2prime-asintuxf3tica-afinada-1}

Una vez extraído el término lineal explícito \(n/6\), el error restante
queda acotado: se tiene \(M(n)=n^2/6+n/6+\theta_n\), con
\(\theta_n\in(-1,\,1/24]\). Esto se sigue de inmediato de
\(24\,M(n)\le(2n+1)^2<24\,M(n)+24\) y \((2n+1)^2=4n^2+4n+1\). Está
formalizado para \(n:\mathbb Z\), sin hipótesis de cota inferior, como
\nolinkurl{PaperII.phiTau_max_sandwich} (huella axiomática
\nolinkurl{propext}, \nolinkurl{Classical.choice}, \nolinkurl{Quot.sound}).

\subsubsection{Observación (coherencia con la cota de Paper
I)}\label{observaciuxf3n-coherencia-con-la-cota-de-paper-i}

donde el lado derecho es la estimación finita
\(|E(G)|-2\nu_3^*(G)\le n^2/6+n/2\) de Paper I {[}4{]}. Esta comparación
es numérica y no una consecuencia de la inclusión de clases, pues la
clase cordal es mayor que la clase split. La holgura es
\(\bigl(n^2/6+n/2\bigr)-M(n)=n/3+O(1)\). La desigualdad está formalizada
como \nolinkurl{PaperII.phiTau_max_le_paperI_bound} (huella axiomática
\nolinkurl{propext}, \nolinkurl{Quot.sound}).

\subsubsection{Definición (defecto de
copia)}\label{definiciuxf3n-defecto-de-copia}

Para \(u,v\) no adyacentes, definamos \[
\begin{aligned}
\Delta_{\Phi_\tau}(u,v)
&=\Phi_\tau(G_{v\to u})+\Phi_\tau(G_{u\to v})-2\Phi_\tau(G)\ge0
\quad(\text{Lema 3.1}),\\
\Gamma_{\Phi_\tau}(u,v)
&=\max\{\Phi_\tau(G_{v\to u}),\Phi_\tau(G_{u\to v})\}-\Phi_\tau(G).
\end{aligned}
\] Entonces \(\Gamma_{\Phi_\tau}\ge\Delta_{\Phi_\tau}/2\), formalizado
como \nolinkurl{copyDefect_nonneg} y
\nolinkurl{copyGamma_ge_half_copyDefect}. La cantidad
\(\Delta_{\Phi_\tau}\) es el recurso que deberá controlar cualquier
enunciado futuro de estabilidad; aquí se registra como una herramienta
cuantitativa, no como un teorema de estabilidad.

\subsubsection{Corolario 7.2 (conjuntos de nivel del funcional de
cobertura)}\label{corolario-7.2-conjuntos-de-nivel-del-funcional-de-cobertura}

Escribamos \(M=2n+1\). Para \(p\) en la rama saturada \(q\ge p-1\) y
todo \(\delta\ge0\), \[
\Phi_\tau(S_{p,q})\ \ge\ \frac{M^2}{24}-\delta\qquad\Longleftrightarrow\qquad \Big|p-\frac{M}{6}\Big|\ \le\ \sqrt{\frac{2\delta}{3}},
\] donde \(\delta\) se mide desde el máximo continuo \(M^2/24\). Es una
equivalencia exacta, obtenida de la identidad (7.1) por monotonía del
cuadrado, y está formalizada como \nolinkurl{level_set_iff}. En
particular, \(\Phi_\tau\) no es plana en ninguna dirección de \(p\), y
sus conjuntos de nivel son intervalos de radio \(\sqrt{2\delta/3}\);
dentro de la familia complete-split, esto resuelve la segunda mitad de
la pregunta de estabilidad planteada en Paper III, §12.4.

\subsection{Referencias}\label{referencias}

{[}1{]} P. Erdős, E. T. Ordman, and Y. Zalcstein, ``Clique partitions of
chordal graphs,'' \emph{Combinatorics, Probability and Computing}
\textbf{2} (1993), no. 4, 409--415.

{[}2{]} M. C. Golumbic, \emph{Algorithmic Graph Theory and Perfect
Graphs}, second edition, Annals of Discrete Mathematics 57, Elsevier,
Amsterdam, 2004.

{[}3{]} J. R. S. Blair and B. Peyton, ``An introduction to chordal
graphs and clique trees,'' in \emph{Graph Theory and Sparse Matrix
Computation}, IMA Volumes in Mathematics and its Applications 56,
Springer, New York, 1993, pp.~1--29.

{[}4{]} J. P. Traverso Gianini, ``Reducción afín de perfiles para
empaquetamientos fraccionales de triángulos en grafos split'' (Paper I
de la serie), preprint, julio de 2026.

{[}5{]} L. de Moura and S. Ullrich, ``The Lean 4 Theorem Prover and
Programming Language,'' in \emph{Automated Deduction -- CADE 28},
Lecture Notes in Computer Science \textbf{12699}, Springer, 2021,
pp.~625--635, \url{https://doi.org/10.1007/978-3-030-79876-5_37}.

{[}6{]} The mathlib Community, ``The Lean Mathematical Library,'' in
\emph{Proceedings of the 9th ACM SIGPLAN International Conference on
Certified Programs and Proofs (CPP 2020)}, ACM, 2020, pp.~367--381,
\url{https://doi.org/10.1145/3372885.3373824}.

{[}7{]} J. P. Traverso Gianini, \emph{Erdős \#81 Chordal Clique
Partitions: Public Preprints and Formalization Artifacts}, public
project repository, 2026,
\url{https://github.com/jtraverso/erdos-81-chordal-clique-partitions},
accessed July 7, 2026.

{[}8{]} T. F. Bloom, ``Erdős Problem \#81,'' \emph{Erdős Problems},
\url{https://www.erdosproblems.com/81}, accessed July 6, 2026.

\end{document}
